\documentclass[11pt]{article}
\usepackage[margin=1in]{geometry}
\usepackage{times}
\usepackage{hyperref}
\hypersetup{colorlinks=true, linkcolor=black, urlcolor=blue, citecolor=black}
\title{A Self-Assessment: Lightweight Skeletons, Mercury Batteries, Solar Slingshots, and Magnetic Floors}
\author{Flyxion \\ \small Independent Researcher}
\date{}
\begin{document}
\maketitle

\begin{abstract}
This paper closes the present series with a self-directed audit rather than a survey of a third party's claim, applying the same evidentiary standard used throughout to four proposals circulated informally within the author's own broader speculative program: replacing human bone with lightweight metal incorporating hydrogen or helium gas cells to reduce body weight; establishing a battery manufacturing base on the premise of abundant mercury supply; using a heat-shielded close solar approach as a gravitational slingshot to reach the outer planets and other stars more quickly; and propelling a craft's internal compartments using electromagnets repelled by magnets embedded in the floor. Each proposal is evaluated on its own physical merits, with credit given where the underlying physics is genuinely sound and clear identification of where it is not.
\end{abstract}

\section{Introduction}

A series built around auditing other people's antigravity, over-unity, and free-energy proposals owes its own speculative program the same treatment, and four proposals recur in that program with enough regularity to warrant the same evidentiary standard applied to Podkletnov, the EmDrive, and the rest of this series: a lightweight metal skeleton incorporating buoyant gas cells, a mercury-based battery manufacturing base, a close solar approach used as a gravitational slingshot toward the outer solar system and beyond, and an internal magnetic propulsion scheme using floor-embedded magnets. They are addressed here in ascending order of how well they survive the audit.

\section{The Magnetic Floor Propulsion Scheme}

The proposal that a craft's internal compartments, or the craft itself, could be propelled by electromagnets aboard the craft repelling against magnets embedded in the craft's own floor is the proposal in this set that most directly reproduces the closed-system momentum objection addressed at length elsewhere in this series regarding the Dean Drive and the Perendev motor: if the floor magnets are part of the same craft as the propelled object, the interaction is entirely internal, and Newton's third law guarantees the forces on the craft and on the propelled mass are equal and opposite, summing to zero net motion of the craft's own center of mass regardless of how the magnets are arranged or timed. The scheme does have a genuine and useful application, however, once its scope is corrected: moving objects or personnel within a stationary or already-moving craft, a maglev-style internal transport system, is an entirely conventional and well precedented application of electromagnetic repulsion, and the proposal's magnetic-floor mechanism is better understood as a design for internal cargo or crew movement than as a propulsion source for the craft as a whole.

\section{The Mercury Battery Factory}

The premise that a large, available supply of mercury could anchor a battery manufacturing base runs into two separate and independent problems rather than one. Mercury oxide batteries are a real, historically used battery chemistry, but they were phased out across most consumer and industrial applications specifically because of mercury's severe toxicity and environmental persistence, a concern formalized internationally in the 2013 Minamata Convention on Mercury, which restricts mercury use in batteries and numerous other products across the great majority of signatory nations, making a new mercury battery manufacturing base a proposal running directly against an actively enforced regulatory trend rather than an underexploited opportunity. Separately, and regardless of the regulatory question, mercury battery chemistry offers energy density substantially below contemporary lithium-based chemistries, meaning even a hypothetical abundant and unregulated mercury supply would anchor a product line already outcompeted on the core performance metric that determines battery selection for nearly every modern application. The proposal's underlying premise, that abundance of a raw material is sufficient grounds for building an industry around it, does not engage with either the safety regulation or the competitive performance question that actually determines a battery chemistry's viability.

\section{The Solar Oberth Slingshot}

Of the four proposals, this one deserves the most credit: using a close, heat-shielded pass by the Sun to gain speed via the Oberth effect, in which a propulsive burn (or, in gravitational-assist framings, the deep gravity well itself) delivers more effective delta-v the faster the craft is already moving and the deeper the gravitational well at the point of the maneuver, is a genuine and actively studied concept in mission design, sometimes called a solar Oberth maneuver, and is closely related to the extreme close solar approach already flown by the Parker Solar Probe, whose own heat shield technology demonstrates that surviving a close perihelion pass is an engineering problem with a real, if extremely demanding, existing solution rather than a purely speculative one. Where the proposal overreaches is in its claimed destination range: a solar Oberth maneuver meaningfully accelerates trajectories toward the outer solar system, and serious mission concepts using exactly this technique for outer planet and near-interstellar-medium missions exist in the aerospace literature, but the resulting velocities, while a substantial improvement over conventional trajectories, remain many orders of magnitude short of what reaching another star system in anything resembling a human timescale would require, meaning the maneuver is a genuine and valuable tool for outer-solar-system and near-interstellar-precursor missions rather than a solution to interstellar travel itself, a distinction the "and other stars much quicker" framing does not maintain.

\section{The Lightweight Metal Skeleton With Gas Cells}

This is the proposal with the least physical viability of the four, on both the materials and the buoyancy side of the claim independently. Bone is not merely a structural strut; it is a living, metabolically active tissue that houses bone marrow, the site of blood cell production, remodels itself continuously in response to mechanical load through a well characterized cellular process, and integrates with surrounding muscle, tendon, and vasculature in ways a metal replacement would need to separately and enormously re-engineer, well beyond the scope of current or foreseeable orthopedic implant technology, which remains limited to replacing specific joints and short bone segments rather than the skeletal system as a whole. Independent of this biological problem, the buoyancy claim does not survive a magnitude check: lifting a person's body weight in air, as opposed to water, requires a genuinely large gas volume, on the order of a cubic meter of helium to offset roughly a kilogram of weight net of the balloon's own material, meaning the comparatively small internal volume available inside a skeleton's marrow cavities, even if entirely evacuated and refilled with helium or hydrogen, would offset a negligible fraction of body weight, several orders of magnitude short of anything a person would perceive as reduced effective weight, while simultaneously eliminating the bone marrow's blood-producing function and introducing a serious decompression and pressure-change hazard at any altitude or depth change, a combination in which the proposal's central goal, a lighter body, is not achieved at any physically meaningful scale even before its biological cost is considered.

\section{Conclusion}

Of the four proposals audited here, the solar Oberth slingshot is grounded in genuine, actively pursued aerospace mission design and needs only a corrected sense of destination scale; the magnetic floor scheme survives as a conventional internal transport mechanism once separated from any claim to propel the craft itself; the mercury battery proposal fails on both regulatory and competitive performance grounds independent of raw material availability; and the lightweight bone and gas cell proposal fails a basic magnitude check on its own stated goal before its biological cost is even weighed, making it the one proposal in this set that this audit cannot recover any functional version of.

\end{document}
